\subsection{The \texorpdfstring{$T$}{T} recursion and the \texorpdfstring{$4k+1$}{4k+1} family}
\label{sec:splittable-4k+1}
This section defines the recursive pair $T_{k,2^l}$, proves its exact cost,
and proves the coefficient-triangular remainder lemma on which all later uses
of $T$ rely.  It then adds five outer pivots to obtain the splittable
degree-$(4k+1)$ family.  \Cref{fig:T-tower} pictures the tower of known
powers that the recursion consumes.

\begin{figure}[H]
  \centering
  \resizebox{\textwidth}{!}{% The tower of "known powers" behind T_{k,2^l}: each rung is one product per
% branch and turns H_{2^l} into a monic H_{2^{l+1}} = H_{2^l}^2 + lower-degree
% corrections carrying fresh parameters; k halves at every rung.
\begin{tikzpicture}[>=Latex, font=\small,
    pow/.style={fp active, rounded corners, minimum width=22mm,
      minimum height=7mm, font=\small},
    gad/.style={fp second, rounded corners, font=\footnotesize,
      align=center, inner sep=2pt}]
  \node[pow] (H2) at (0,0) {$H_2 = x^2+\dots$};
  \node[pow] (H4) at (0,1.7) {$H_4 = H_2^2+\dots$};
  \node[pow] (H8) at (0,3.4) {$H_8 = H_4^2+\dots$};
  \node[pow] (H16) at (0,5.1) {$H_{16} = H_8^2+\dots$};
  \node at (0,6.1) {$\vdots$};

  \draw[fp flow] (H2) -- node[right, xshift=1mm, font=\footnotesize]
    {$(H_2+S_1)(H_2-S_1)+S_2$} (H4);
  \draw[fp flow] (H4) -- node[right, xshift=1mm, font=\footnotesize]
    {$(H_4+S_1)(H_4-S_1)+S_2$} (H8);
  \draw[fp flow] (H8) -- node[right, xshift=1mm, font=\footnotesize]
    {$\bigl((H_8+S_1)+S_2\bigr)\bigl((H_8+S_1)-S_2\bigr)+S_3$} (H16);

  % gadgets feeding the corrections (left column)
  \node[gad, left=12mm of H4] (g4) {$S_1 = x+a,\; S_2=e$\\ (even step, $l{=}1$)};
  \node[gad, left=12mm of H8] (g8) {$S_1 = H_2+Q_1,\; S_2 = Q_1$\\ (even step, $l{=}2$)};
  \node[gad, left=12mm of H16] (g16) {$S_1=H_4+Q_3,\; S_2=H_2+Q_1,\; S_3=Q_1$\\ (odd step, $l{=}3$)};
  \draw[fp dependence] (g4) -- ($(H2)!0.5!(H4)$);
  \draw[fp dependence] (g8) -- ($(H4)!0.5!(H8)$);
  \draw[fp dependence] (g16) -- ($(H8)!0.5!(H16)$);

  \node[font=\footnotesize, text=black!85, align=left, text width=118mm,
    anchor=north west] at (-4.6,-0.7)
    {One product per branch per rung; only the two identified bases share
     products.  The parameters inside $S_1,S_2,S_3$ are fresh and sit
     strictly below the leading term, so the top of $T$ equals $H^{\,k}$ and
     the corrections are read off top-down.};
\end{tikzpicture}
}
  \caption{The tower of known powers used by the recursion $T_{k,2^l}$ (\Cref{alg:constr-Tk2l,alg:constr-Tk2l-base}).
  In each branch, a rung uses one product
  $(H+S_1)(H-S_1)+S_2$ (even step) or
  $((H+S_1)+S_2)((H+S_1)-S_2)+S_3$ (odd step), producing a monic
  $H_{2^{l+1}}=H_{2^l}^2+\text{lower terms}$; those lower terms carry fresh
  parameters through small $Q$-gadgets.  The index $k$ halves at every rung:
  at level $l$ the recursion computes
  $T_{k,2^l}=H_{2^l}^{\,k}+(\text{terms of degree}\le 2^l(k-1))$, and the
  base $T_{1,2^L}=H_{2^L}$ costs no product.
  The two branches normally use parallel products.  Product sharing occurs
  only in the explicitly identified $l=1$ even base and the immediately
  following odd $l=2$ base, where the required scalar-difference invariant
  has already been established.}
  \label{fig:T-tower}
\end{figure}

We use named parameter blocks in the recursive definition.  If $\mathsf B$
is a block of $q$ fresh parameters, $Q_q[\mathsf B]$ denotes the corresponding
known-powers gadget, with its required power arguments suppressed.  A display
$[\mathsf B_1\mid\cdots\mid\mathsf B_s]$ assigns
$\alpha_0,\ldots,\alpha_{d-1}$ to the blocks from left to right.  The sizes
shown below sum to $d=(k-1)2^l$, so this notation is only a readable
reindexing of one consecutive parameter list.  Recursive $T$ calls inherit
$x$ and the existing power tower; their displays show only the fresh block
and the newly formed top power(s).

\begin{algorithm}[H]
    \caption{Recursive construction of $T_{k,D}$, ordinary branches
    ($D=2^l$; $D\ge4$ in the even branch and $D\ge8$ in the odd branch).}
    \label{alg:constr-Tk2l}
    \begin{algorithmic}
        \State Put $H=H_D$ and $\widetilde H=\widetilde H_D$.
        \If{$k=1$}
            \State $T^{(1)}_{1,D}=H$, \qquad $T^{(2)}_{1,D}=\widetilde H$.
        \ElsIf{$k=2m$}
            \State Put $r=D/2$, $b=(k-2)D$, and allocate
            $[\mathsf I_b\mid\sigma\mid\mathsf B^-_{r-1}
              \mid\delta\mid\mathsf B^+_{r-1}]$.
            \State $U_1=H_r+Q_{r-1}[\mathsf B^+]$,
            $V_1=Q_{r-1}[\mathsf B^-]$,
            $U_2=H_r+\delta$, $V_2=\sigma$.
            \State $H_{2D}=H^2-U_1^2+V_1$, \qquad
            $\widetilde H_{2D}=\widetilde H^2-U_2^2+V_2$.
            \State $T^{(1)}_{k,D}=T^{(1)}_{m,2D}[\mathsf I](H_{2D})$,
            \qquad
            $T^{(2)}_{k,D}=T^{(2)}_{m,2D}[\mathsf I]
              (H_{2D},\widetilde H_{2D})$.
        \Else \Comment{$k=2m+1$}
            \State Put $r=D/2$, $q=D/4$, $b=(k-2)D$, and allocate
            \Statex \hspace{\algorithmicindent}$[\alpha_0\mid\mathsf B^{\rm low}_{D-1}
              \mid\mathsf I_{b-D}\mid\zeta\mid\mathsf B^-_{q-1}
              \mid\varepsilon\mid\mathsf B^0_{q-1}
              \mid\delta\mid\mathsf B^+_{r-1}]$.
            \State $U_1=H_r+Q_{r-1}[\mathsf B^+]$,
            $V_1=H_q+Q_{q-1}[\mathsf B^0]$,
            $W_1=Q_{q-1}[\mathsf B^-]$.
            \State $U_2=H_r+\delta$, $V_2=H_q+\varepsilon$,
            $W_2=\zeta$.
            \State $H_{2D}=(H+U_1)^2-V_1^2+W_1$, \qquad
            $\widetilde H_{2D}=(\widetilde H+U_2)^2-V_2^2+W_2$.
            \State $T^{(1)}_{k,D}=(H-(k-1)U_1)
              T^{(1)}_{m,2D}[\mathsf I](H_{2D})
              +Q_{D-1}[\mathsf B^{\rm low}]$.
            \State $T^{(2)}_{k,D}=(\widetilde H-(k-1)U_2)
              T^{(2)}_{m,2D}[\mathsf I](H_{2D},\widetilde H_{2D})
              +\alpha_0$.
        \EndIf
    \end{algorithmic}
\end{algorithm}

The two shared-product bases are invoked only when
$\widetilde H_D-H_D$ is scalar.  The even $D=2$ base preserves that
difference at degree four, which is exactly the hypothesis needed by the
following odd $D=4$ base.

\begin{algorithm}[H]
    \caption{Shared-product bases for $T_{k,D}$ ($D=2$ with $k$ even,
    and $D=4$ with $k$ odd).}
    \label{alg:constr-Tk2l-base}
      \begin{algorithmic}
          \If{$k=2m$ and $D=2$}
              \State Allocate $[\mathsf I_{2k-4}\mid e\mid a]$ and set
              $H_4=H_2^2-(x+a)^2+e$.
              \State Retain the supplied scalar shift
              $\rho=\widetilde H_2-H_2$ and set
              $\widetilde H_4=H_4+\rho$.
              \Comment{do not recompute the difference}
              \State $T^{(1)}_{k,2}=T^{(1)}_{m,4}[\mathsf I](H_4)$,
              \qquad
              $T^{(2)}_{k,2}=T^{(2)}_{m,4}[\mathsf I]
                 (H_4,\widetilde H_4)$.
          \Else \Comment{$k=2m+1\ge3$ and $D=4$}
              \State Allocate
              $[\alpha_0\mid\mathsf B^{\rm low}_3
                \mid\mathsf I_{4(k-3)}\mid z\mid w\mid v\mid u]$.
              \State $U_1=H_2+x+u$, $V_1=x+v$, $W_1=w$, and
              $H_8=(H_4+U_1)^2-V_1^2+W_1$.
              \State Retain $\rho=\widetilde H_4-H_4$.  For the proof put
              $U_2=U_1-\rho$, $V_2=V_1$, $W_2=W_1+z$, but compute only
              $\widetilde H_8=H_8+z$.
              \State Put $F_1=H_4-(k-1)U_1$ and use the shared affine form
              $F_2=F_1+k\rho
                    =\widetilde H_4-(k-1)U_2$.
              \State $T^{(1)}_{k,4}=F_1
                T^{(1)}_{m,8}[\mathsf I](H_8)
                +Q_3[\mathsf B^{\rm low}]$.
              \State $T^{(2)}_{k,4}=F_2
                T^{(2)}_{m,8}[\mathsf I](H_8,\widetilde H_8)
                +\alpha_0$.
          \EndIf
      \end{algorithmic}
  \end{algorithm}
\begin{remark}[Why the shared odd base is admissible]
    The condition that $\tilde H_4-H_4$ is scalar says exactly that the two
    quartics have the same nonconstant coefficients.  The difference may be
    an active parameter; once $H_4$ and $\tilde H_4$ are supplied to the
    recursive call, that scalar is recoverable from their constant
    coefficients and belongs to the given data.  It is not an additional
    condition on the target polynomial.  The shared odd $l=2$ branch is invoked
    only immediately after the even $l=1$ branch, which constructs
    \[
       \tilde H_4-H_4=\tilde H_2-H_2.
    \]
    Thus the recursion establishes the condition before it uses it.  A
    standalone odd $T_{k,4}$ call with arbitrary quartic inputs is not covered
    by this shared-base definition (it would require a separate two-product
    variant).  No such call occurs in the exact-count construction.
\end{remark}

\begin{lemma}[Exact cost of the $T$ recursion]\label{lem:T-multiplication-count}
    Let $\mu(k,l)$ be the number of multiplication gates used to compute the
    pair $(T^{(1)}_{k,2^l},T^{(2)}_{k,2^l})$ when
    $(x,H_2,\ldots,H_{2^l},\tilde H_{2^l})$ is already available.  For every
    call occurring in the construction---including the shared $l=1$ step and,
    when reached from it, the shared odd $l=2$ step---we have
    \[
       \mu(k,l)=(k-1)2^{l-1}.
    \]
\end{lemma}
\begin{proof}
    Additions, scalar shifts, and multiplication by the displayed integers are
    free in this count; the latter can be implemented by repeated additions.
    The case $k=1$ uses no multiplication and agrees with the formula.

    Suppose first that $k$ is even.  If $l=1$, the shared base computes $H_4$
    with one product and obtains $\tilde H_4=H_4+(\tilde H_2-H_2)$ without a
    second product.  If $l\ge2$, each of the two
    $Q_{2^{l-1}-1}$ blocks costs $2^{l-2}-1$ multiplications by
    \Cref{lem:fill-Q-count}; for $l=2$ these are $Q_1=x+\alpha_0$,
    using zero products in agreement with the formula.  The two new
    degree-$2^{l+1}$ powers cost one product each.  Thus in both cases the even recurrence is
    \[
       \mu(k,l)=\mu(k/2,l+1)+2^{l-1}.
    \]

    Now suppose that $k\ge3$ is odd.  In the shared $l=2$ base there is one
    product for the common octic core, one for $Q_3$, and two products for the
    final recursive factors.  The shift $\tilde H_8=H_8+z$ costs
    nothing, so the overhead is $4=2^l$.  For $l\ge3$, the four embedded
    Mersenne gadgets have costs
    \[
       (2^{l-2}-1),\quad (2^{l-3}-1),\quad
       (2^{l-3}-1),\quad (2^{l-1}-1).
    \]
    The two new powers and the two final recursive factors cost four more
    products.  Their sum is exactly $2^l$.  Hence the odd recurrence is
    \[
       \mu(k,l)=\mu((k-1)/2,l+1)+2^l.
    \]

    Strong induction on $k$ now closes the count.  In the even case,
    \[
      \mu(k,l)=\left(\frac{k}{2}-1\right)2^l+2^{l-1}
               =(k-1)2^{l-1},
    \]
    while in the odd case,
    \[
      \mu(k,l)=\left(\frac{k-1}{2}-1\right)2^l+2^l
               =(k-1)2^{l-1}.
    \]
\end{proof}

\begin{lemma}[Cubic degree loss in an odd $T$ step]
\label{lem:odd-T-cubic-loss}
Let $k\ge3$, let $D=2r$ with $r\ge1$, and let $H,U$ be monic of
degrees $D,r$.  Put
\[
  c_k=\binom{k}{2},\qquad \tau_k=2\binom{k}{3}
      =\frac{k(k-1)(k-2)}3,
\]
and
\[
  E(H,U)=\bigl(H-(k-1)U\bigr)(H+U)^{k-1}-H^k
          +c_kU^2H^{k-2}.
\]
Then
\[
  \deg E(H,U)\le (k-2)D+r,
  \qquad
  [x^{(k-2)D+r}]E(H,U)=-\tau_k.
\]
\end{lemma}
\begin{proof}
The binomial theorem gives
\[
  \bigl(H-(k-1)U\bigr)(H+U)^{k-1}-H^k
  =\sum_{t=1}^{k}
    \left(\binom{k-1}{t}-(k-1)\binom{k-1}{t-1}\right)
    H^{k-t}U^t,
\]
where $\binom{k-1}{k}=0$.  The coefficients for $t=1,2,3$ are
$0,-c_k,-\tau_k$.  The added quadratic term cancels $t=2$.
The $t=3$ term has degree $(k-3)D+3r=(k-2)D+r$, while every
$t\ge4$ term has degree $kD-tr\le(k-2)D$.  Monicity of $H,U$
therefore gives both assertions.
\end{proof}

\begin{lemma}[Boundary transport of a remainder pair]
\label{lem:seam-transport}
Let $A_1,A_2$ be monic of the same degree $f\ge0$ and put
$\varphi=[x^{f-1}]A_1$ (so $\varphi=0$ when $f=0$).  Let $e\ge0$ and let
$B_1,B_2$ have degree at most $e$ with
\[
  [x^{e}]B_1=[x^{e}]B_2=\ell,\qquad [x^{e-1}]B_1=a,
\]
where a coefficient of negative index is read as zero.  Then the shifted
combination $P=xA_1B_1+A_2B_2$ has degree at most $f+e+1$, and its two top
rows are
\begin{equation}
  [x^{f+e+1}]P=\ell,\qquad
  [x^{f+e}]P=a+(\varphi+1)\ell.
  \tag{seam-transport}\label{eq:seam-transport}
\end{equation}
In particular both values are determined by the interface
$(\ell,a,\varphi)$ alone: they involve neither $[x^{e-1}]B_2$, nor any
lower coefficient of $B_1,B_2$, nor any coefficient of $A_2$ below the
leading one.
\end{lemma}
\begin{proof}
Both products have degree at most $f+e$, so $\deg P\le f+e+1$,
\[
  [x^{f+e+1}]P=[x^{f+e}](A_1B_1),
  \qquad
  [x^{f+e}]P=[x^{f+e-1}](A_1B_1)+[x^{f+e}](A_2B_2).
\]
If $e\ge1$, the top-two specialization \eqref{eq:monic-top-two} of
\Cref{lem:monic-cauchy-transport} gives $[x^{f+e}](A_iB_i)=[x^e]B_i=\ell$
for $i=1,2$ and
$[x^{f+e-1}](A_1B_1)=[x^{e-1}]B_1+[x^{f-1}]A_1\,[x^{e}]B_1=a+\varphi\ell$.
If $e=0$, then $B_1=B_2=\ell$ are scalars, $a=0$, and the same three
values are read off directly from $A_i\ell$.  Adding the contributions
gives \eqref{eq:seam-transport}; the final sentence records which
coefficients the two formulas use.
\end{proof}

\begin{lemma}\label{lem:Rk2l}
    Let $l\ge 2$ and $k\ge 1$, and supply monic polynomials
    $H_2,\ldots,H_{2^l},\tilde H_{2^l}$ of their indicated degrees.  In the
    exceptional case where $l=2$ and $k\ge3$ is odd, assume that
    $\tilde H_4-H_4$ is a scalar and use the shared-product base in
    \Cref{alg:constr-Tk2l-base}.
    Assume that $\mathbb F$ is $(k2^l)$-admissible.
    Define the remainder polynomials by
    \[
      R^{(1)}_{k,2^l}\coloneqq T^{(1)}_{k,2^l}(x,H_2,\ldots,H_{2^l})-H_{2^l}^k,
      \qquad
      R^{(2)}_{k,2^l}\coloneqq T^{(2)}_{k,2^l}(x,H_2,\ldots,H_{2^l},\tilde H_{2^l})-\tilde H_{2^l}^k.
    \]
    Then:
    \begin{enumerate}
        \item By definition,
        \(
          T^{(1)}_{k,2^l}=H_{2^l}^k+R^{(1)}_{k,2^l}
        \)
        and
        \(
          T^{(2)}_{k,2^l}=\tilde H_{2^l}^k+R^{(2)}_{k,2^l}.
        \)
        \item If $k=1$, both remainders vanish.  If $k\ge2$, then
        $\deg R^{(1)}_{k,2^l}=\deg R^{(2)}_{k,2^l}=(k-1)2^l$.
        \item With $d=(k-1)2^l$, the remainder pair is
        coefficient-triangular in the sense of
        \Cref{def:coefficient-triangular}, in the row order specified by
        the stage tables below (with each known-powers block ordered as in
        \Cref{lem:Q-unitriangular}).  Hence the raw parameters
        $\alpha_0,\ldots,\alpha_{d-1}$ are extractable from
        $xR^{(1)}_{k,2^l}+R^{(2)}_{k,2^l}$, and
        $(T^{(1)}_{k,2^l},T^{(2)}_{k,2^l})$ is compatible on $\rng d$,
        given $(H_2,\ldots,H_{2^l},\tilde H_{2^l})$.
        \item If $k\ge2$, put $D=2^l$ and retain the monic
        degree-$D/2$ auxiliary polynomials $S^{(i)}_1$ from the defining
        branch.  With $H^{(1)}=H_D$, $H^{(2)}=\widetilde H_D$,
        $\sigma_i=[x^{D/2-1}]S^{(i)}_1$, and
        $h_i=[x^{D-1}]H^{(i)}$, one has
        \[
          [x^d]R^{(i)}_{k,D}=-\gamma_k,
          \qquad
          [x^{d-1}]R^{(i)}_{k,D}
            =-\gamma_k\bigl(2\sigma_i+(k-2)h_i\bigr),
        \]
        where
        $\gamma_k=k/2$ for even $k$ and
        $\gamma_k=k(k-1)/2$ for odd $k$.
    \end{enumerate}
\end{lemma}
\begin{figure}[H]
  \centering
  \resizebox{\textwidth}{!}{% Decoder-stage map for the three branches of lem:Rk2l.
\begin{tikzpicture}[x=1cm,y=1cm,>=Latex]
  \draw[fp decode] (0,6.15) -- node[fp label, above] {descending coefficient degree}
    (13.65,6.15);
  \node[fp label, anchor=east] at (0,6.15) {high};
  \node[fp label, anchor=west] at (13.65,6.15) {low};

  % Even branch.
  \node[fp panel title] at (0,5.65) {even $k=2m$};
  \node[fp active, minimum width=21mm, anchor=west] (ep) at (0,5.02)
    {$Q_+$\\$\lambda=-k$};
  \node[fp pivot, minimum width=9mm, anchor=west] (ed) at (2.1,5.02)
    {$\delta$\\$-k$};
  \node[fp second, minimum width=20mm, anchor=west] (em) at (3.0,5.02)
    {$Q_-$\\$m$};
  \node[fp pivot, minimum width=9mm, anchor=west] (es) at (5.0,5.02)
    {$\sigma$\\$m$};
  \node[fp recursive, minimum width=48mm, anchor=west] (er) at (5.9,5.02)
    {recursive block $R'_{m,2D}$\\(no degree shift)};
  \draw[fp dependence] (es.south east) to[bend right=12]
    node[fp label, below] {tail known} (er.south west);

  % Shared odd base.
  \node[fp panel title] at (0,3.87) {odd $k$, shared base $D=4$};
  \node[fp active, minimum width=13mm, anchor=west] (ou) at (0,3.22)
    {$u$\\$-k(k-1)$};
  \node[fp second, minimum width=12mm, anchor=west] (ov) at (1.3,3.22)
    {$v$\\$-(k-1)$};
  \node[fp pivot, minimum width=10mm, anchor=west] (ow) at (2.5,3.22)
    {$w$\\$m$};
  \node[fp pivot, minimum width=10mm, anchor=west] (oz) at (3.5,3.22)
    {$z$\\$m$};
  \node[fp recursive, minimum width=39mm, anchor=west] (orr) at (4.5,3.22)
    {inner pivots shifted by $4$\\through $L_1,L_2$};
  \node[fp second, minimum width=22mm, anchor=west] (oq) at (8.4,3.22)
    {low $Q_3$\\unit pivots};
  \node[fp pivot, minimum width=8mm, anchor=west] at (10.6,3.22)
    {$\alpha_0$\\$1$};
  \draw[fp dependence] (oz.south east) to[bend right=12]
    node[fp label, below] {form shared powers} (orr.south west);
  \node[fp label, anchor=west, text=fpGray] at (11.65,3.22)
    {$u,v,w,z$ share\\the outer product};

  % Ordinary odd branch.
  \node[fp panel title] at (0,2.05) {odd $k$, ordinary branch $D\ge8$};
  \node[fp active, minimum width=16mm, anchor=west] (qp) at (0,1.38)
    {$Q_+$\\$-k(k-1)$};
  \node[fp pivot, minimum width=15mm, anchor=west] (delta) at (1.6,1.38)
    {$\delta$\\$-k(k-1)$};
  \node[fp second, minimum width=14mm, anchor=west] (qz) at (3.1,1.38)
    {$Q_0$\\$-(k-1)$};
  \node[fp pivot, minimum width=13mm, anchor=west] (eps) at (4.5,1.38)
    {$\varepsilon$\\$-(k-1)$};
  \node[fp second, minimum width=12mm, anchor=west] (qm) at (5.8,1.38)
    {$Q_-$\\$m$};
  \node[fp pivot, minimum width=8mm, anchor=west] (zeta) at (7.0,1.38)
    {$\zeta$};
  \node[fp recursive, minimum width=28mm, anchor=west] (rr) at (7.8,1.38)
    {inner recursion\\shifted by $D$};
  \node[fp second, minimum width=18mm, anchor=west] (lq) at (10.6,1.38)
    {low $Q_{D-1}$\\unit pivots};
  \node[fp pivot, minimum width=8mm, anchor=west] at (12.4,1.38)
    {$\alpha_0$\\$1$};
  \node[fp seam, minimum width=7mm, minimum height=3mm] at (2.4,0.62) {};
  \node[fp label, anchor=west] at (2.82,0.62)
    {named boundary corrections, subtracted at the junctions
     tail$\to$recursion and recursion$\to$low block};
  \draw[fp dependence] (zeta.south east) to[bend right=10] (rr.south west);
  \node[fp label, anchor=west, text=fpBlue!80!black] at (0,0.12)
    {tail pivots $\longrightarrow$ shifted recursive pivots $\longrightarrow$ low unitriangular pivots};
\end{tikzpicture}
}
  \caption{Decoder map for the three branches of the central
  $T_{k,2^l}$ remainder lemma.  Each coloured segment is a consecutive row
  interval; its label is the active parameter block and its constant pivot
  slope.  In every branch the fresh tail is decoded first.  The resulting
  powers and monic factors then expose the recursive block, after which the
  low unitriangular $Q$-block is read in the two odd branches; the even
  branch has no low block.  Boundary corrections within and between these
  segments are listed in the stage tables and the formulas below.}
  \label{fig:Rk2l-stages}
\end{figure}
\begin{proof}
    \emph{Reading guide.}
    The algebra below certifies the three coloured rows of
    \Cref{fig:Rk2l-stages}.  Every branch follows the same four-step template:
    derive one exact remainder identity; identify its highest nonzero term;
    verify the tail pivots and their two seam rows; then transport the recursive
    block and, in the odd branches, append the low unitriangular block.
    On a first reading, one may
    follow the displayed identity, stage table, and final concatenation in each
    branch.  The paragraphs headed ``seam rows'' give the boundary
    corrections within and between the blocks in the figure.
    \Cref{alg:decode-Rk2l} restates the same
    argument as an executable decoder.

    Assertion (1) is the definition.  We prove (2--4) simultaneously by
    induction on $k$.  Carrying the two coefficients in (4) is essential:
    they are the parameter-free corrections at the boundaries between two
    consecutive decoder blocks.
    For $k=1$, we have $T^{(1)}_{1,2^l}=H_{2^l}$ and $T^{(2)}_{1,2^l}=\tilde H_{2^l}$ by the defining algorithm, hence $R^{(1)}_{1,2^l}=R^{(2)}_{1,2^l}=0$ and the claims are trivial.

    We may therefore assume $k\ge2$ below.  For the remainder of the proof
    put
    \[
       D=2^l,\qquad d=(k-1)D.
    \]
    In an even step the recursive call has size
    $(k/2)(2D)=kD$, and in an odd step it has size
    $((k-1)/2)(2D)=(k-1)D<kD$.  Thus $kD$-admissibility supplies the
    induction hypothesis in both cases.
    In the odd branch with $l=2$ this automatically means $k\ge3$, so all
    indices in the shared-base formulas are defined.

      Within each known-powers block, parameter pivots use the recursive
      row order of \Cref{lem:Q-unitriangular}, rather than the raw order in
      \Cref{alg:constr-known-2n-1}.  The stage tables first recover factor
      coefficients; the known-powers decoder then recovers the raw block.
      Row permutations are confined to their displayed blocks and do not
      change any coefficient cutoff.  In the proof and stage tables below,
      $\alpha_j$ denotes the parameter assigned to row $j$ by this blockwise
      permutation; the construction and final decoder retain the raw order.

      Each branch is now proved in the same order: one exact remainder
      identity, its degree/support bounds, its tail stage table, and finally
      the recursive and low blocks.  Every product transport below is an
      instance of the causal monic Cauchy rule
      \eqref{eq:monic-Cauchy}; together with
      \Cref{lem:Q-unitriangular}, it says that the current row has unit slope
      in the current factor coefficient and otherwise uses only higher rows.
      Every seam correction is an instance of \Cref{lem:seam-transport}:
      the two rows at which a decoder block ends receive the top two
      coefficients of the block below, carried across the monic factor that
      multiplies it and across the shift $x$ of the combined polynomial.
      The inner remainder pair is carried this way in all three branches,
      with the interface $(\ell,a,\varphi)$ fixed once and for all in the
      next paragraph; the even branch carries one further pair, its
      quadratic binomial term.

      \smallskip\noindent\emph{Inner boundary values.}
      Put $m=\lfloor k/2\rfloor$ and $h=[x^{D-1}]H_D$, recall
      \[
        \gamma_t=\begin{cases}t/2&t\text{ even},\\ t(t-1)/2&t\text{ odd},\end{cases}
        \qquad \gamma_1=0,
      \]
      and define the three fixed seam values
      \[
        \ell_m=-\gamma_m,\qquad
        a_m=-2\gamma_m\bigl(1+(m-1)h\bigr),\qquad
        \theta_m=a_m+(h+1)\ell_m+\mathbf 1_{\{k=3\}}.
        \tag{inner-boundary}\label{eq:inner-boundary}
      \]
      Let $R'^{(i)}=R^{(i)}_{m,2D}$ be the inner remainder pair of the
      recursive call at $(m,l+1)$ and put $e=(m-1)2D$.  We claim
      \[
        [x^{e}]R'^{(1)}=[x^{e}]R'^{(2)}=\ell_m,
        \qquad
        [x^{e-1}]R'^{(1)}=a_m .
        \tag{inner-top-two}\label{eq:inner-top-two}
      \]
      If $m=1$ both inner remainders vanish and $\gamma_1=0$, so there is
      nothing to prove.  If $m\ge2$, assertion~(4) of the induction
      hypothesis at $(m,l+1)$ gives $[x^e]R'^{(i)}=-\gamma_m$ and
      $[x^{e-1}]R'^{(1)}=-\gamma_m\bigl(2\sigma'_1+(m-2)h'_1\bigr)$, where
      $\sigma'_1=[x^{D-1}]S'^{(1)}_1$ and $h'_1=[x^{2D-1}]H_{2D}$ refer to
      the inner call.  Since $l+1\ge3$, that call is an ordinary branch of
      \Cref{alg:constr-Tk2l}, so $S'^{(1)}_1=H_D+Q_{D-1}[\mathsf B^+]$ with
      $Q_{D-1}$ monic by \Cref{lem:Q-unitriangular}, whence
      $\sigma'_1=h+1$.  In each of the three branches below, $H_{2D}$ is
      $H_D^2$ or $(H_D+U_1)^2$ plus terms of degree at most $D$, where
      $\deg U_1\le D/2<D-1$; hence $h'_1=2h$.  Substituting,
      $[x^{e-1}]R'^{(1)}=-\gamma_m\bigl(2(h+1)+2(m-2)h\bigr)=a_m$.
      The second inner branch has the same leading coefficient $\ell_m$;
      its next coefficient is never needed, because by
      \Cref{lem:seam-transport} it does not reach a seam row.

      In all three tables, $K$ consists only of field constants and the
      supplied power coefficients (including the tilded top power).  We use
      the stage-table convention above with
      $\alpha_{>j}=(\alpha_{j+1},\ldots,\alpha_{d-1})$.  In particular,
      quantities decoded in a higher row are substituted; they are not added
      to $K$ as free side information.

      \paragraph{Even $k$.}
      Let $k=2m$, put $b=(k-2)D$ and $r=D/2$, and write
      \[
        U=S^{(1)}_1=H_r+Q_+,\quad V=S^{(1)}_2=Q_-,\quad
        \widetilde U=S^{(2)}_1=H_r+\delta,\quad
        \widetilde V=S^{(2)}_2=\sigma.
      \]
      Here $Q_+$ and $Q_-$ have degree $r-1$, while $\delta$ and
      $\sigma$ are scalars.  In the formulas below $H^{(1)}=H_D$ and
      $H^{(2)}=\widetilde H_D$, and

      \[
      E_1=-U^2+V,\qquad E_2=-\widetilde U^2+\sigma,\qquad
      R'_i=R^{(i)}_{m,2D}.
      \]
      \smallskip\noindent\emph{Identity, degree, and top boundary.}
      The tail parameters have the order
      \[
        (\sigma,\eta^-_0,\ldots,\eta^-_{r-2},\delta,
          \eta^+_0,\ldots,\eta^+_{r-2})
        =(\alpha_{b},\ldots,\alpha_{b+D-1}),
        \tag{R-even-block}
      \]
      where the $\eta$'s are the row-ordered pivot parameters of the two
      $Q$-polynomials.
      The recursive definition and the binomial theorem give the single
      identity
      \[
        R^{(i)}_{k,D}=mE_i(H^{(i)})^{k-2}+\widehat R_i+R'_i,
        \tag{R-even-exp}
      \]
      with $\deg\widehat R_i,\deg R'_i\le b$.  Explicitly,
      $\widehat R_i=\sum_{q=2}^{m}\binom mq E_i^q(H^{(i)})^{k-2q}$;
      terms with $q\ge3$ have degree at most $b-D$.

      The summand $-m(S^{(i)}_1)^2(H^{(i)})^{k-2}$ has degree $d$ and
      leading coefficient $-m$.  The $S^{(i)}_2$ companion has degree at
      most $b+r-1$, while $\widehat R_i$ and $R'_i$ have degree at most
      $b$.  Hence the remainder has exact degree $d$, and its top two
      coefficients are
      \begin{equation}
        [x^d]R^{(i)}_{k,D}=-m,\qquad
        [x^{d-1}]R^{(i)}_{k,D}
          =-m\bigl(2\sigma_i+(k-2)h_i\bigr),
        \tag{R-top-two-even}\label{R-top-two-even}
      \end{equation}
      where $\sigma_i=[x^{r-1}]S^{(i)}_1$ and
      $h_i=[x^{D-1}]H^{(i)}$.  This proves assertions~(2) and~(4) for the
      even branch.

      \smallskip\noindent\emph{Tail pivots and support.}
      Applying \eqref{eq:monic-Cauchy} to the first-order terms gives the complete tail
      table
      \[
      \begin{array}{c|c|c|c}
        \text{rows }j&\text{active block}&\lambda_j&\text{source}\\ \hline
        b+D-1,\ldots,b+r+1
          &\alpha_{b+r+1},\ldots,\alpha_{b+D-1}&-2m=-k&-xU^2(H^{(1)})^{k-2}\\
        b+r&\alpha_{b+r}=\delta&-2m=-k&-\widetilde U^2(H^{(2)})^{k-2}\\
        b+r-1,\ldots,b+1
          &\alpha_{b+1},\ldots,\alpha_{b+r-1}&m&xV(H^{(1)})^{k-2}\\
        b&\alpha_b=\sigma&m&\sigma(H^{(2)})^{k-2}\\
        b-1,\ldots,0&\text{recursive block}&\text{inner pivot}&R'_1,R'_2
      \end{array}
      \tag{R-even-table}
      \]
      For example, the pivot parameter $\eta^+_s$ first occurs as
      $2\eta^+_s$ plus higher-row terms in $U^2$ at degree $r+s$; the $x$-shift and the monic
      factor $(H^{(1)})^{k-2}$ put it in row $b+r+s+1$, with slope $-2m$.
      The other rows follow from the same Cauchy-product calculation.  In a
      first-branch row $j$ the displayed occurrence lies in
      $[x^{j-1}]R^{(1)}$, and replacing any monic leading term by a lower
      coefficient can only lower its degree; hence that parameter occurs in
      the first component only through degree $j-1$.  The analogous
      second-branch occurrence lies in $[x^j]R^{(2)}$ and nowhere above
      degree $j$.  This proves the two component-support inequalities for all
      four tail blocks, not just the combined pivot formulas.  At the
      two bottom tail rows one must retain the boundary error

      \[
        E^{\rm bd}=x\widehat R_1+\widehat R_2+xR'_1+R'_2.
      \]
      \smallskip\noindent\emph{Seam rows.}
      Let $s=[x^{r-1}]U$; here $h=[x^{D-1}]H^{(1)}$ is the value in
      \eqref{eq:inner-boundary}.  The two seam rows $b+1$ and $b$ of
      $E^{\rm bd}$ come from two instances of \Cref{lem:seam-transport},
      both landing on these rows because $f+e=b$ in each:
      \begin{itemize}
        \item the quadratic binomial term
        $\binom m2E_i^2(H^{(i)})^{k-4}$ (for $m=1$ the sum $\widehat R_i$ is
        empty and $\binom m2=0$), with
        $(A_i,B_i,f,e)=\bigl((H^{(i)})^{k-4},\binom m2E_i^2,(k-4)D,2D\bigr)$:
        here $\varphi=(k-4)h$; $\ell=\binom m2$ because
        $E_i^2=U_i^4+\cdots$ is monic of degree $2D$; and
        $a=\binom m2\,[x^{2D-1}]U^4=4s\binom m2$, because $U^2V$ has degree
        at most $3r-1<2D-1$;
        \item the recursive pair, with $(A_i,B_i,f,e)=(1,R'_i,0,b)$: here
        $\varphi=0$ and $(\ell,a)=(\ell_m,a_m)$ by \eqref{eq:inner-top-two},
        since $e=(m-1)2D=b$ in the even branch.
      \end{itemize}
      The $q\ge3$ terms of $\widehat R_i$ have degree at most $b-D<b-1$ and
      meet neither row.  Adding the two instances of
      \eqref{eq:seam-transport} gives
      \[
        [x^{b+1}]E^{\rm bd}=\binom m2-\gamma_m,\qquad
        [x^{b}]E^{\rm bd}
        =\binom m2\bigl(4s+(k-4)h+1\bigr)+a_m+\ell_m,
        \tag{R-even-boundary-values}\label{eq:R-even-boundary-values}
      \]
      where $\ell_m=-\gamma_m$ and
      $a_m=-2\gamma_m(1+(m-1)h)$ as in \eqref{eq:inner-boundary} (with both
      zero for $m=1$).
      These are constants once the already exposed tail is substituted.
      \Cref{lem:seam-transport} also shows explicitly why no inner parameter
      is present in either boundary equation: the recursive pair enters only
      through its interface $(\ell_m,a_m,\varphi)$, and each of these values
      is fixed by \eqref{eq:inner-top-two}.
      Subtracting them gives the pivots in the last two tail rows.  Every
      other summand in a displayed row either has lower degree or contains a
      strictly higher coefficient of the relevant $Q$-polynomial.  This proves
      the support and pivot assertions in the tail.

      \smallskip\noindent\emph{Recursive block and concatenation.}
      Once this tail is known, the new powers $H_{2D}$ and
      $\widetilde H_{2D}$ are known as well.  The remaining parameters are the
      unshifted inner block $\alpha_0,\ldots,\alpha_{b-1}$.  Since the even
      recursion is the direct pair $(R'_1,R'_2)$ plus known tail terms, the
      induction hypothesis applies at the same rows $0,\ldots,b-1$ (there is
      no degree shift in the even branch).  The side data
      $H_{2D},\widetilde H_{2D}$ and the coefficients of the other displayed
      terms are polynomially recoverable from
      $(K;\alpha_b,\ldots,\alpha_{d-1})$.  Consequently the tail rows and the
      inner rows are glued by
      \Cref{lem:triangular-block-concatenation}; in particular, treating the
      newly formed powers as known here is a substitution of higher-block
      formulas, not an enlargement of the cutoff.  Finally the top-two coefficient
      calculation in \eqref{R-top-two-even} shows that all four coefficients
      in degrees $d$ and $d-1$ are parameter-free (the top coefficients of
      the displayed $S_1$ blocks are fixed by monicity and the known powers).

      \paragraph{Odd $k$, shared $l=2$ base.}
      Put $D=4$, $d_0=4(k-2)$, $d=4(k-1)$, $m=(k-1)/2$, and
      $\rho=\widetilde H_4-H_4$; since both quartics are supplied,
      $\rho\polyfrom K$ and is an $x$-independent known scalar.
      The base definitions give
      \[
      \begin{aligned}
        S^{(1)}_1&=H_2+x+u,&
        S^{(1)}_2&=x+v,&
        S^{(1)}_3&=w,\\
        S^{(2)}_1&=S^{(1)}_1-\rho,&
        S^{(2)}_2&=S^{(1)}_2,&
        S^{(2)}_3&=w+z,
      \end{aligned}
      \]
      with $(z,w,v,u)=(\alpha_{d_0},\alpha_{d_0+1},
      \alpha_{d_0+2},\alpha_{d_0+3})$.  The shared identities
      $\widetilde H_4+S^{(2)}_1=H_4+S^{(1)}_1$ and
      $\widetilde H_8=H_8+z$ show that both branches use the same product.
      Here $h=[x^3]H_4$, and $\ell_m,a_m,\theta_m$ are the inner boundary
      values \eqref{eq:inner-boundary}; put
      \[
        \tau=\frac{k(k-1)(k-2)}3 .
      \]
      \smallskip\noindent\emph{Tail pivots, seam rows, and top boundary.}
      The factors $F_1,F_2$ of \Cref{alg:constr-Tk2l-base} are monic of
      degree $4$ with $[x^3]F_1=h$, because $\deg U_1=2$.  The inner pair
      at $(m,3)$ has degree at most $e=8(m-1)=d_0-4$ and the top two
      coefficients \eqref{eq:inner-top-two}.  \Cref{lem:seam-transport}
      with $(A_1,A_2,B_1,B_2,f,e)=(F_1,F_2,R'^{(1)},R'^{(2)},4,d_0-4)$ and
      $\varphi=h$ shows that the product $xF_1R'^{(1)}+F_2R'^{(2)}$
      contributes exactly $\ell_m$ in the $w$ row $d_0+1$ and
      $a_m+(h+1)\ell_m$ in the $z$ row $d_0$, and nothing above.  The low
      term $xQ_3$ contributes its known leading coefficient $1$ in row $4$,
      which is the $z$ row exactly when $k=3$; this is the indicator in
      $\theta_m$.
      After subtracting the terms depending on already exposed parameters,
      the first row therefore has the fixed correction $-m-\tau$, the $w$
      row has the fixed recursive correction $\ell_m$, and the $z$ row has
      the fixed correction $\theta_m$.  Direct expansion of the four
      boundary rows then gives
      \[
      \begin{array}{c|c|c}
        \text{row }j&\text{parameter}&\lambda_j\\ \hline
        d-1&u&-k(k-1)\\
        d-2&v&-(k-1)\\
        d-3&w&m\\
        d-4=d_0&z&m
      \end{array}
      \tag{R-odd-base-table}
      \]
      after these corrections and the preceding rows have been substituted.
      To verify the table symbolically, inspect the highest row in
      $D_R=xR^{(1)}+R^{(2)}$ in which each scalar can occur.  The scalar $u$
      is the constant coefficient of the monic quadratic $S^{(1)}_1$ (and of
      $S^{(2)}_1$ up to the known shift $\rho$); in the principal term
      $-\binom{k}{2}(S^{(1)}_1)^2H_4^{k-2}$ its first-branch occurrence is
      shifted by $x$ to row $d-1$, with slope
      $-2\binom{k}{2}=-k(k-1)$.  Every occurrence of $u$ in the second branch
      or in a term containing at least three copies of $S_1$ is at least one
      row lower.  Similarly, $v$ first occurs through the cross term in
      $-m(S^{(1)}_2)^2$, in row $d-2$ with slope $-2m$; $w$ first occurs in
      the first-branch linear $mS^{(1)}_3$ term, in row $d-3$ with slope $m$;
      and $z$ first occurs in the corresponding second-branch term, in row
      $d-4$ with slope $m$.  The monic factors multiplying these terms can
      contribute only higher already-known coefficients by
      \eqref{eq:monic-Cauchy}.
      This degree inspection also proves the component-support condition: in
      the first component
      a parameter assigned to row $j$ occurs only through degree $j-1$, and
      in the second only through degree $j$.  Thus the four displayed slopes
      are constant and nonzero under the stated admissibility hypothesis.
      The same four-row expansion also proves the strengthened boundary
      assertion: with $\sigma_i=[x^{1}]S^{(i)}_1$ and
      $h_i=[x^3]H^{(i)}$, where $H^{(1)}=H_4$ and
      $H^{(2)}=\widetilde H_4$,
      \[
        [x^d]R^{(i)}_{k,4}=-\frac{k(k-1)}2,\qquad
        [x^{d-1}]R^{(i)}_{k,4}
          =-\frac{k(k-1)}2\bigl(2\sigma_i+(k-2)h_i\bigr).
      \]
      These are parameter-free because the constant shifts in the displayed
      $S_1$ blocks do not change $\sigma_i$.
      The first coefficient is nonzero by admissibility, and therefore also
      proves $\deg R^{(i)}_{k,4}=d$ in the shared-base branch.
      \smallskip\noindent\emph{Shifted recursive and low blocks.}
      After the tail is recovered, the inner pair is multiplied by the monic
      degree-four factors
      $H_4-(k-1)S^{(1)}_1$ and
      $\widetilde H_4-(k-1)S^{(2)}_1$.  At row $4$, the term
      $[x^4](xQ_3)=1$ is a known boundary contribution.  If $k>3$ it is
      subtracted before the first shifted pivot in
      \Cref{lem:triangular-shift}; if $k=3$ the shifted interval is empty,
      and the same contribution is the final tail boundary already recorded by the
      $\mathbf 1_{\{k=3\}}$ term in $\theta_m$.  The shifted induction
      hypothesis therefore applies to the rows $4,\ldots,d_0-1$, including
      the overlap of the two factors.  The coefficients of those factors are
      polynomially recoverable from
      $(K;\alpha_{d_0},\ldots,\alpha_{d-1})$, so
      \Cref{lem:triangular-shift} and
      \Cref{lem:triangular-block-concatenation} justify the substitution of
      the decoded tail.  The low rows are the unitriangular $Q_3$ block and
      the scalar $\alpha_0$; \Cref{lem:Q-unitriangular} supplies the pivots on
      rows $1,2,3,0$.  This proves the shared-base case without a
      separate composition argument.

      \paragraph{Odd $k$, $l\ge3$.}
      Put $r=D/2$, $q=D/4$, $b=(k-2)D$, $m=(k-1)/2$, and
      $c=k(k-1)/2$.
      Abbreviate
      \[
        U=H_r+Q_+,\quad V=H_q+Q_0,\quad W=Q_-,
        \qquad
        \widetilde U=H_r+\delta,\quad
        \widetilde V=H_q+\varepsilon,\quad
        \widetilde W=\zeta.
      \]
      The tail order is
      \[
        (\zeta,Q_-\text{-block},\varepsilon,Q_0\text{-block},
          \delta,Q_+\text{-block})
        =(\alpha_{b},\ldots,\alpha_{b+D-1}),
        \tag{R-odd-block}
      \]
      with block lengths $1,q-1,1,q-1,1,r-1$.  Set
      \[
      \begin{aligned}
        L_1&=H-(k-1)U,&
        L_2&=\widetilde H-(k-1)\widetilde U,\\
        K_1&=L_1(H+U)^{k-3},&
        K_2&=L_2(\widetilde H+\widetilde U)^{k-3},\\
        G_1&=-V^2+W,&
        G_2&=-\widetilde V^2+\zeta.
      \end{aligned}
      \]
      \smallskip\noindent\emph{Identity and error separation.}
      The factors $L_i$ are monic of degree $D$ and the factors $K_i$ are
      monic of degree $b$.  An exact binomial expansion
      of the odd recursion gives
      \begin{align}
        R^{(1)}_{k,D}&=-cU^2H^{k-2}+mG_1K_1
             +L_1R'^{(1)}+Q_{D-1}+E_1,\\
        R^{(2)}_{k,D}&=-c\widetilde U^2\widetilde H^{k-2}+mG_2K_2
             +L_2R'^{(2)}+\alpha_0+E_2,
        \tag{R-odd-block-exp}
      \end{align}
      where $R'$ is the inner remainder pair at $(m,l+1)$.  We split the
      error, since a single coarse degree bound does not prove the required
      support at the two adjacent windows.  With
      $(H^{(1)},U_1)=(H,U)$ and $(H^{(2)},U_2)=(\widetilde H,\widetilde U)$, set
      \[
      \begin{aligned}
      E_i^{U}&\coloneqq
        L_i(H^{(i)}+U_i)^{k-1}-(H^{(i)})^k+cU_i^2(H^{(i)})^{k-2},\\
      E_i^{G}&\coloneqq
        L_i\sum_{q=2}^{m}\binom mqG_i^q
                       (H^{(i)}+U_i)^{k-1-2q}.
      \end{aligned}
      \]
      (An empty sum is zero.)  Then $E_i=E_i^U+E_i^G$.  Put
      $\tau=2\binom{k}{3}=k(k-1)(k-2)/3$.
      \Cref{lem:odd-T-cubic-loss}, applied to $(H^{(i)},U_i)$, gives the first
      two bounds below.  The $q$th summand of $E_i^G$ has degree at most

      \[
        D+qr+(k-1-2q)D
          =kD-\frac{3qD}{2}\le b-D\qquad(q\ge2),
      \]
      and therefore
      \[
        \deg E_i^U\le b+r=d-r,\qquad
        [x^{b+r}]E_i^U=-\tau,\qquad
        \deg E_i^G\le b-D.
        \tag{R-odd-error-bound}\label{eq:R-odd-error-bound}
      \]
      The coefficient at degree $b+r$ is parameter-free.  Hence $xE_1^U$
      meets the tail first at row $b+r+1$
      and $E_2^U$ first at row $b+r$, with a fixed correction in either
      case; every lower coefficient of $E_i^U$ depends only on the
      $Q_+$/$\delta$ block, which has already been exposed when that row is
      processed.  More precisely, using the coefficient of $x^s$ in $Q_+$
      in a term of $E_1^U$ lowers its degree by at least $r-s$, so its
      occurrence in $xE_1^U$ is in a row at most $b+s+1$, strictly below its
      pivot row $b+r+s+1$.  Using the scalar $\delta$ instead of the leading
      term of $\widetilde U$ lowers the degree by $r$, so every occurrence of
      $\delta$ in $E_2^U$ has degree at most $b$, below its pivot row $b+r$.
      Finally, the last bound in \eqref{eq:R-odd-error-bound} puts every
      $G$-error below the entire tail window.  These are the precise
      degree-loss statements used in the stage table.

      \smallskip\noindent\emph{Degree and top boundary.}
      In particular, the principal term
      $-cU_i^2(H^{(i)})^{k-2}$ has degree $d$ and leading coefficient $-c$;
      the $mG_iK_i$ and $E_i^U$ terms have degree at most $d-r$, the
      product $L_iR'^{(i)}$ has degree at most $b$, and the remaining terms
      are lower still.  Since $c$ is invertible under the admissibility
      hypothesis, this proves the exact-degree assertion in the ordinary odd
      branch.

      The same degree estimate gives the other part of the strengthened
      induction invariant.  The first displayed correction is the only term
      that can reach degrees $d$ and $d-1$; its top two coefficients are
      $-1$ and $-(2\sigma_i+(k-2)h_i)$, where
      $\sigma_i=[x^{r-1}]S^{(i)}_1$ and
      $h_i=[x^{D-1}]H^{(i)}$.  Consequently
      \begin{equation}
        [x^d]R^{(i)}_{k,D}=-c,\qquad
        [x^{d-1}]R^{(i)}_{k,D}
          =-c\bigl(2\sigma_i+(k-2)h_i\bigr).
        \tag{R-top-two-odd}\label{R-top-two-odd}
      \end{equation}
      For the shared $D=4$ base the same statement follows from the four-row
      expansion above (the error terms have degree at most $d-2$).

      \smallskip\noindent\emph{Tail pivots and support.}
      The resulting stage table is
      \[
      \begin{array}{c|c|c|c}
        \text{rows }j&\text{active block}&\lambda_j&\text{source}\\ \hline
        d-1,\ldots,b+r+2&Q_+\text{-block}&-2c=-k(k-1)&-cxU^2H^{k-2}\\
        b+r+1&\text{lowest $Q_+$ coefficient}&-2c=-k(k-1)&-cxU^2H^{k-2}\\
        b+r&\delta&-2c=-k(k-1)&-c\widetilde U^2\widetilde H^{k-2}\\
        b+r-1,\ldots,b+q+1&Q_0\text{-block}&-2m=-(k-1)&mxG_1K_1\\
        b+q&\varepsilon&-2m=-(k-1)&mG_2K_2\\
        b+q-1,\ldots,b+1&Q_-\text{-block}&m&mxWK_1\\
        b&\zeta&m&mG_2K_2
      \end{array}
      \tag{R-odd-table}\label{eq:R-odd-table}
      \]
      The source column refers to the combined polynomial
      $D_R=xR^{(1)}+R^{(2)}$.  Here is the support check, separated by
      source.  In the first source, \eqref{eq:monic-Cauchy} isolates the coefficient of
      $U$ (respectively $\widetilde U$) at the current row and all other
      summands use a strictly higher coefficient.  Since $U$ has the
      coefficient-triangular $Q_+$ block and $\widetilde U$ has only the
      scalar $\delta$, this gives the first three rows of the table.  The
      term $E^U_i$ has the same dependence: its top coefficient is the fixed
      $-\tau$, and replacing a leading coefficient of $U_i$ by its scalar
      term loses $r$ degrees, so $\delta$ cannot occur before the $\delta$
      row.  The term $E^G_i$ is below $b-D$ and hence cannot meet any tail
      row.  In the second source, \eqref{eq:monic-Cauchy} applied to $G_iK_i$ first
      exposes the coefficient of $V$ or $W$ at the current row; the
      coefficients of $K_i$ and all remaining summands depend only on the
      already exposed $Q_+$ and $\delta$ blocks.  This gives the four
      stage--2 rows and their slopes.  As in the even case, a first-branch
      pivot in row $j$ is an occurrence in $[x^{j-1}]R^{(1)}$ and no occurrence
      of that active parameter can lie higher, while a second-branch pivot is
      an occurrence in $[x^j]R^{(2)}$ and none can lie higher.  The degree
      bounds for $E^U,E^G$ just proved cover the only terms not visible in
      those two Cauchy-product calculations.  Thus, after substitution, every
      row in the table has exactly the stage-table support form, with the
      displayed nonzero pivot.

      \smallskip\noindent\emph{Seam rows.}
      At row $b+r+1$ the stage-$2$ leading term and
      the cubic term in $E_1$ contribute the fixed correction
      $-m-\tau$, where $\tau=k(k-1)(k-2)/3$.  At rows $b+1$ and $b$ the
      inner pair is carried across $L_1,L_2$.  These factors are monic of
      degree $D$ with $[x^{D-1}]L_1=h$, since $\deg U=r<D-1$; the inner
      pair has degree at most $e=(m-1)2D=b-D$ and the top two coefficients
      \eqref{eq:inner-top-two}.  \Cref{lem:seam-transport} with
      $(A_1,A_2,B_1,B_2,f,e)=(L_1,L_2,R'^{(1)},R'^{(2)},D,b-D)$ and
      $\varphi=h$ therefore shows that $xL_1R'^{(1)}+L_2R'^{(2)}$
      contributes $\ell_m$ at row $b+1$ and $a_m+(h+1)\ell_m$ at row $b$,
      and nothing above row $b+1$ (when $m=1$, equivalently $k=3$, both
      inner remainders vanish and $\ell_m=a_m=0$).  The latter row also
      receives the known leading coefficient of $xQ_{D-1}$ exactly when
      $k=3$, giving $\theta_m$ as in \eqref{eq:inner-boundary}.  A
      corresponding fixed contribution from the same
      stage--$2$/cubic terms can occur in the $\delta$ row $b+r$; it is
      absorbed into the known correction term $F$ there and does not change
      the pivot slope.  After these corrections and the values of the
      already exposed higher blocks have been subtracted, every row in the
      table is exactly the indicated affine pivot.  In particular, the scalar
      $\varepsilon$ is solved before the $Q_-$ ($S_3$) window.

      \smallskip\noindent\emph{Shifted recursive and low blocks.}
      After the tail is known, apply \Cref{lem:triangular-shift} with shift
      $D$ to $(L_1R'^{(1)},L_2R'^{(2)})$ and the induction hypothesis.  The
      inner block has length $b-D$, and every coefficient of $L_i$ is
      polynomially recoverable from
      $(K;\alpha_b,\ldots,\alpha_{d-1})$.  The higher-block clause in
      \Cref{lem:triangular-shift} therefore moves the inner pivots, including
      the overlap of the two factors, to rows $D,\ldots,b-1$ without changing
      their cutoffs.  The principal and error terms outside the recursive
      product now depend only on the decoded tail block, so they are the
      permitted higher-block additions in \Cref{lem:triangular-shift} and are
      subtracted row by row.  At row $D$, $xQ_{D-1}$ contributes its known leading
      coefficient $1$, which is subtracted before the first shifted pivot
      (when $k=3$, this is the boundary row $b=D$ already accounted for).
      Below row $D$, the remaining active block is
      $Q_{D-1}[\alpha_1,\ldots,\alpha_{D-1}]$ together with $\alpha_0$;
      \Cref{lem:Q-unitriangular} gives unit pivots on rows
      $1,\ldots,D-1,0$.  The block-concatenation lemma now glues the tail,
      shifted inner block, and low block, proving the support and pivot
      conditions
      for every row: the three intervals are respectively
      $b,\ldots,d-1$, $D,\ldots,b-1$, and $0,\ldots,D-1$
      (with an empty middle interval when $k=3$).  The top-boundary assertion
      follows from the simultaneous
      two-term coefficient calculation above.  This completes the odd case.

      In each of the even, shared-base, and ordinary odd branches, the glued
      stage table is precisely a coefficient-triangular certificate.
      Applying \Cref{lem:triangular-implies-compatible} gives the extraction
      and causal compatibility claimed in~(3); the simultaneous top-two
      calculation in the same branch gives~(4).  This completes the
      induction and the proof of (2--4).

\end{proof}

\begin{algorithm}[H]
  \caption{Decoder summary for $xR^{(1)}_{k,2^l}+R^{(2)}_{k,2^l}$}\label{alg:decode-Rk2l}
  \begin{algorithmic}
    \Require integers $l\ge 2$, $k\ge 1$, known powers $(H_2,\ldots,H_{2^l})$ and the shifted power $\tilde H_{2^l}$, and $P=xR^{(1)}_{k,2^l}+R^{(2)}_{k,2^l}$; if $l=2$ and $k\ge3$ is odd, require $\tilde H_4-H_4$ to be a scalar (possibly zero)
    \Ensure the parameter block $\alpha_0,\ldots,\alpha_{(k-1)2^l-1}$ (and the derived inputs for the recursive call at $(k',l+1)$, $k'\in\{k/2,(k-1)/2\}$)

    \If{$k=1$}
      \State Output the empty parameter list.
    \ElsIf{$k$ is even}
      \State Use the top-degree window to peel the monic factor in the first-order correction term, recovering $S^{(1)}_1$ (via \Cref{lem:peel-monic-factor} and \Cref{lem:monic-from-power} with $m=2$).
      \State Recover the scalar shift in $S^{(2)}_1=H_{2^{l-1}}+\delta$ from the boundary coefficient of the $(S^{(2)}_1)^2$ branch using \Cref{lem:scalar-shift-square}.
      \State Recover the coefficients of $S^{(1)}_2$ in descending order; the descent ends at degree $(k-2)2^l+1$, where the parameter-free boundary correction in \eqref{eq:R-even-boundary-values} is subtracted.  Decode its embedded $Q$ block.
      \State Recover the scalar $S^{(2)}_2$ from the following boundary row (it enters affinely with slope $\tfrac{k}{2}$).
      \State Derive $H_{2^{l+1}}$ and $\tilde H_{2^{l+1}}$ from the recovered $(S^{(1)}_1,S^{(1)}_2,S^{(2)}_1,S^{(2)}_2)$.
      \State Isolate the recursive combined remainder polynomial $xR^{(1)}_{k/2,2^{l+1}}+R^{(2)}_{k/2,2^{l+1}}$ and recurse at parameters $(k/2,l+1)$.
    \ElsIf{$k\ge3$ is odd and $l=2$}
      \State Recover $u,v,w,z$ from the four descending affine pivots in the shared-base part of the proof of \Cref{lem:Rk2l}.
      \State Derive $H_8$, $\tilde H_8=H_8+z$, and the two degree-$4$ factors multiplying the recursive pair.
      \State Apply the shift-$4$ triangular recurrence to the overlapping monic-factor products, invoking the inner decoder at $((k-1)/2,3)$; then subtract the recovered inner contribution and decode the low $Q_3\mathbin{\oplus}\alpha_0$ block.
    \Else \Comment{$k$ odd and $l\ge3$; $Q_\pm,Q_0$ are the $Q_{r-1}[\mathsf B^\pm]$, $Q_{q-1}[\mathsf B^0]$ blocks of \Cref{alg:constr-Tk2l}}
      \State Read the $Q_+$ block and then the scalar $\delta$ from the first three row blocks of \eqref{eq:R-odd-table}, subtracting the fixed correction $-m-\tau$ at their common boundary.
      \State Read the $Q_0$ block and $\varepsilon$, followed by the $Q_-$ block and $\zeta$, from the remaining four row blocks.  At the last two rows subtract the parameter-free inner boundary coefficients $\ell_m$ and $\theta_m$ of \eqref{eq:inner-boundary} (carried across $L_1,L_2$ by \Cref{lem:seam-transport}).
      \State Form the two monic factors $L_1,L_2$.  Descend through the overlapping polynomial $xL_1R'^{(1)}+L_2R'^{(2)}$ by the shifted triangular recurrence of \Cref{lem:triangular-shift}, invoking the inner decoder at $((k-1)/2,l+1)$; do not divide the two overlapping products separately.
      \State Subtract the recovered inner contribution and decode the low $Q_{2^l-1}\mathbin{\oplus}\alpha_0$ block.
    \EndIf
  \end{algorithmic}
\end{algorithm}

\begin{lemma}[Leading coefficients of the remainder polynomials]\label{lem:Rk2l-leading-coeff}
    Assume the hypotheses in \Cref{lem:Rk2l}, and let $k\ge2$ and
    $D=2^l$.  Put $d=(k-1)D$ and
    \[
      \gamma_k\coloneqq
      \begin{cases}k/2&k\text{ even},\\ k(k-1)/2&k\text{ odd},\end{cases}
      \qquad
      \sigma_i\coloneqq [x^{D/2-1}]S^{(i)}_1,\quad
      h_i\coloneqq [x^{D-1}]H^{(i)},
    \]
    where $H^{(1)}=H_D$ and $H^{(2)}=\widetilde H_D$.  Then, for $i=1,2$,
    \[
      [x^d]R^{(i)}_{k,D}=-\gamma_k,
      \qquad
      [x^{d-1}]R^{(i)}_{k,D}
        =-\gamma_k\bigl(2\sigma_i+(k-2)h_i\bigr).
      \tag{R-top-two}
    \]
    In particular, the leading-coefficient formula used below is the first
    equality.
\end{lemma}
\begin{proof}
    This is assertion~(4) of \Cref{lem:Rk2l}, recorded separately for later
    reference.  Its simultaneous induction proof is the top-two coefficient
    calculation in the even, shared-base, and ordinary odd branches there.
\end{proof}

\begin{lemma}[Parameter-free top boundary]\label{lem:Rk2l-top-boundary}
    Under the hypotheses of \Cref{lem:Rk2l}, assume $k\ge2$ and put
    $d=(k-1)2^l$.
    All four coefficients
    \[
      [x^d]R^{(1)}_{k,2^l},\quad [x^{d-1}]R^{(1)}_{k,2^l},\quad
      [x^d]R^{(2)}_{k,2^l},\quad [x^{d-1}]R^{(2)}_{k,2^l}
    \]
    are derivable from the given-power data alone; in particular, they are
    independent of the parameter block of the $T_{k,2^l}$ call.
\end{lemma}
\begin{proof}
    By \Cref{lem:Rk2l-leading-coeff}, these coefficients are determined by
    the top two coefficients of $H^{(i)}$ and $S^{(i)}_1$.  In every even-$k$
    call (including $l=2$), and in every odd-$k$ call with $l\ge3$,
    $S^{(1)}_1=H_{2^{l-1}}+Q$ with $Q$ monic of degree
    $2^{l-1}-1$, while $S^{(2)}_1=H_{2^{l-1}}+$ a scalar.  Hence their top
    two coefficients are fixed by the given powers and monicity.  In the
    shared $l=2$ base, $S^{(1)}_1=H_2+x+u$ and
    $S^{(2)}_1=S^{(1)}_1-\rho$; again $u$ and $\rho$ affect only the constant
    term.  Substitution in (R-top-two) proves the claim.
\end{proof}

\begin{lemma}[Causal substitution for a perturbed top power]
\label{lem:causal-perturbed-T}
Let $D=2^l\ge4$, put $r=D/2$, and let $M=2k\ge2$ and $N=MD$.
Fix monic lower powers $H_2,\ldots,H_{D/2}$ as given data.  Let $H=H_D$
be a given monic power, let $Q$ be monic of degree $r-1$,
and let $\delta$ be an unknown $x$-independent scalar.  Set
\[
  \widehat H=H+Q,\qquad \widetilde H=\widehat H+\delta .
\]
Let
\[
  S^{(1)}=T^{(1)}_{M,D}(\widehat H),\qquad
  S^{(2)}=T^{(2)}_{M,D}(\widehat H,\widetilde H),\qquad
  \Psi=xS^{(1)}+S^{(2)} .
\]
Assume that $\mathbb F$ is $N$-admissible.  Then
$(S^{(1)},S^{(2)})$ is compatible on $\rng{N-r}$ given the original
known powers, and $\widehat H,\widetilde H$ are polynomially recoverable
from those powers and the corresponding coefficients of $\Psi$.
\end{lemma}
\begin{proof}
Because $M$ is even, this call is not in the exceptional shared odd base.
Thus $N=MD$-admissibility lets us apply
\Cref{lem:Rk2l,lem:Rk2l-top-boundary} at $(M,l)$.
Put $d=N-D$ and write
\[
  Q=x^{r-1}+\sum_{q=0}^{r-2}q_qx^q .
\]
The active quantities occur in the following descending row blocks:
\[
\begin{array}{c|c|c}
  \text{rows of }\Psi&\text{new quantity}&\text{pivot slope}\\ \hline
  d+r-1,\ldots,d+1&q_{r-2},\ldots,q_0&M\\
  d&\delta&M\\
  d-1,\ldots,0&\text{parameters of }R_{M,D}&
       \text{slopes from \Cref{lem:Rk2l}}
\end{array}
\tag{perturbed-T-table}
\]
We justify every row and its cutoff.

First, the binomial expansion and $\deg Q=r-1<D$ give
\[
 [x^{d+q}]\widehat H^M
   =M q_q+F_q(q_{q+1},\ldots,q_{r-2};H)
   \qquad(0\le q\le r-2),
\tag{perturbed-power-pivot}
\]
because a term containing two copies of $Q$ has degree at most $d-2$.
The quantity $q_q$ cannot occur one degree higher, in either
$\widehat H^M$ or $\widetilde H^M$.  The remainders do not meet these rows,
except that $xR^{(1)}$ contributes the fixed coefficient
$[x^d]R^{(1)}=-M/2$ in the last row $d+1$.  Consequently row $d+q+1$ of
$\Psi$ has the first pivot in the table, plus a polynomial in the already
recovered $q_{q+1},\ldots,q_{r-2}$ and the original known powers.  Reading
$q=r-2,r-3,\ldots,0$ therefore recovers all coefficients of $Q$, and hence
$\widehat H$.

At row $d$, a scalar perturbation contributes
\[
 [x^d]\widetilde H^M=[x^d]\widehat H^M+M\delta .
\]
The remaining boundary terms $[x^{d-1}]R^{(1)}$ and
$[x^d]R^{(2)}$ are parameter-free by
\Cref{lem:Rk2l-top-boundary}.  More explicitly, their formulas use only
$[x^{D-1}]\widehat H=[x^{D-1}]\widetilde H=[x^{D-1}]H$ and the fixed top
two coefficients of the auxiliary $S_1$ blocks.  Thus row $d$ has slope
$M$ in $\delta$ and recovers it without circular side information.

After $Q$ and $\delta$ have been recovered, subtract the two powers:
\[
 \Psi-\bigl(x\widehat H^M+\widetilde H^M\bigr)
       =xR^{(1)}+R^{(2)} .
\]
The coefficient-triangular conclusion of \Cref{lem:Rk2l} supplies the last
block of the table.

It remains only to check that this decoder is causal.  A coefficient $q_u$
can occur in either perturbed power only in degrees at most $d+u$; it is
recovered in row $d+u+1$.  The scalar $\delta$ can occur only in the second
component and only in degrees at most $d$; it is recovered in row $d$.
Any occurrence of these quantities in a lower remainder coefficient is
therefore already known before that row is processed.  The last block then
has exactly the two cutoffs asserted by \Cref{lem:Rk2l}: row $j+1$ suffices
for the first component and row $j$ for the second.  Finally, coefficients
in degrees at least $N-r$ agree with the known polynomial $H^M$, while the
coefficient in degree $N-r-1=d+r-1$ is fixed by the monicity of $Q$.
This proves compatibility on $\rng{N-r}$ and also recovers
$\widehat H=H+Q$ and $\widetilde H=\widehat H+\delta$.
\end{proof}

\begin{lemma}[The $4k+1$ family is splittable]\label{lem:4k+1-splittable}
    Let $k\ge1$ and assume that $\mathbb F$ is $(4k+1)$-admissible.
    Define the quadratic base
    \[
      H_2[\alpha_{4k-1},\alpha_{4k}](x)\coloneqq (x+\alpha_{4k})x+\alpha_{4k-1},
      \qquad
      \tilde H_2\coloneqq H_2+\alpha_{4k-2}.
    \]
    Let
    \[
      T^{(1)}_{4k+1}\coloneqq T^{(1)}_{2k,2}[\alpha_0,\ldots,\alpha_{4k-3}](x,H_2),
      \qquad
      T^{(2)}_{4k+1}\coloneqq T^{(2)}_{2k,2}[\alpha_0,\ldots,\alpha_{4k-3}](x,H_2,\tilde H_2),
    \]
    and set $P_{4k+1}\coloneqq xT^{(1)}_{4k+1}+T^{(2)}_{4k+1}$.
    Then $P_{4k+1}$ is decodable, and the pair
    $(T^{(1)}_{4k+1},T^{(2)}_{4k+1})$ is compatible on
    $\rng{4k+1}$ with no auxiliary data.  In particular it is a splittable
    pair for $4k+1$.  Moreover, the displayed construction is a joint
    $2k$-realization and records its monic $H_2$ and $H_4$.
\end{lemma}

\Cref{fig:4k1-crown} pictures the five-pivot crown over the $T_{k,4}$ core
inside the $T_{2k,2}$ call.

\begin{figure}[H]
  \centering
  \resizebox{\textwidth}{!}{% Forward construction and inverse five-row crown for the 4k+1 family.
\begin{tikzpicture}[x=1cm,y=1cm,>=Latex]
  \node[fp panel title] at (0,5.05) {(a) Forward construction};
  \node[fp active, minimum width=28mm] (h2) at (1.55,4.22)
    {$H_2=(x+b)x+c$\\one product};
  \node[fp second, minimum width=37mm] (h4) at (5.2,4.22)
    {$H_4=H_2^2-(x+a)^2+e$\\one shared product};
  \node[fp pivot, minimum width=26mm] (ht) at (8.72,4.22)
    {$\widetilde H_4=H_4+\rho$\\no product};
  \node[fp recursive, minimum width=27mm] (t) at (11.75,4.22)
    {$(U,V)=T_{k,4}$\\recursive core};
  \draw[fp flow] (h2) -- (h4);
  \draw[fp flow] (h4) -- (ht);
  \draw[fp flow] (ht) -- (t);
  \node[fp label, below=2mm of t] {$P_{4k+1}=xU+V$};

  \node[fp panel title] at (0,3.05) {(b) Inverse coefficient crown};
  \draw[fp decode] (0,2.62) -- node[fp label, above]
    {top five rows, then the $T$ remainder} (13.55,2.62);
  \node[fp active, minimum width=15mm, anchor=west] (b) at (0,1.92)
    {$b$\\$2k$};
  \node[fp active, minimum width=15mm, anchor=west] (c) at (1.5,1.92)
    {$c$\\$2k$};
  \node[fp second, minimum width=15mm, anchor=west] (a) at (3.0,1.92)
    {$a$\\$-2k$};
  \node[fp second, minimum width=15mm, anchor=west] (e) at (4.5,1.92)
    {$e$\\$k$};
  \node[fp pivot, minimum width=15mm, anchor=west] (rho) at (6.0,1.92)
    {$\rho$\\$k$};
  \node[fp recursive, minimum width=52mm, anchor=west] (inner) at (7.5,1.92)
    {$x(U-H_4^k)+(V-\widetilde H_4^k)$\\apply the $(R^{(1)}_{k,4},R^{(2)}_{k,4})$ decoder};

  \node[fp known, minimum width=25mm] (dh2) at (1.5,0.72) {reconstruct $H_2$};
  \node[fp known, minimum width=25mm] (dh4) at (4.5,0.72) {reconstruct $H_4$};
  \node[fp known, minimum width=24mm] (dht) at (7.15,0.72) {form $\widetilde H_4$};
  \draw[fp dependence] (b.south) -- (dh2.north west);
  \draw[fp dependence] (c.south) -- (dh2.north east);
  \draw[fp dependence] (a.south) -- (dh4.north west);
  \draw[fp dependence] (e.south) -- (dh4.north east);
  \draw[fp dependence] (rho.south) -- (dht.north);
  \draw[fp dependence] (dh2.east) -- (dh4.west);
  \draw[fp dependence] (dh4.east) -- (dht.west);
  \draw[fp dependence] (dht.east) -- (inner.south west);
  \node[fp label, anchor=west, text=fpBlue!80!black] at (9.05,0.72)
    {subtract the recovered powers\\before entering the inner rows};
\end{tikzpicture}
}
  \caption{The $4k+1$ construction and its inverse crown.  Forward, the
  quadratic and quartic base polynomials are shared by the two branches of the
  $T_{k,4}$ call.  Backward, the top five coefficients expose
  $b,c,a,e,\rho$ with the displayed constant slopes.  These reconstruct
  $H_2,H_4,\widetilde H_4$ before the decoder subtracts their $k$th powers
  and enters the coefficient-triangular remainder block.}
  \label{fig:4k1-crown}
\end{figure}
\begin{proof}
    Put
    \[
      b=\alpha_{4k},\quad c=\alpha_{4k-1},\quad
      \rho=\alpha_{4k-2},\quad a=\alpha_{4k-3},\quad
      e=\alpha_{4k-4}.
    \]
    The shared $l=1$ step is equivalently
    \[
      H_2=x^2+bx+c,\qquad
      H_4=H_2^2-(x+a)^2+e,\qquad
      \tilde H_4=H_4+\rho,
    \]
    followed by the call
    \[
      (U,V)=T_{k,4}[\alpha_0,\ldots,\alpha_{4k-5}]
         (x,H_2,H_4,\tilde H_4),\qquad P_{4k+1}=xU+V.
    \]
    Write $H_4=x^4+h_3x^3+h_2x^2+h_1x+h_0$.  Direct expansion gives
    \[
      h_3=2b,\quad h_2=b^2+2c-1,\quad
      h_1=2bc-2a,\quad h_0=c^2-a^2+e.
    \]

    Here $\tilde H_4-H_4=\rho$ is a scalar, and
    $(4k+1)$-admissibility implies $4k$-admissibility.  Hence all hypotheses
    of \Cref{lem:Rk2l} hold at $(k,l)=(k,2)$, and
    \[
      U=H_4^k+A,\qquad V=(H_4+\rho)^k+B,\qquad
      \deg A,\deg B\le4(k-1).
    \]
    (For $k=1$ both remainders vanish; for $k\ge2$ both inequalities are
    equalities.)
    Hence $xA+B$ has degree at most $4k-3$.  Expanding the top five
    coefficients of $P_{4k+1}$ and substituting previously recovered values
    gives the triangular pivots
    \[
    \begin{array}{c|c|c}
      \text{coefficient}&\text{new parameter}&\text{slope}\\ \hline
      \coeff{P_{4k+1}}{4k}&b&2k\\
      \coeff{P_{4k+1}}{4k-1}&c&2k\\
      \coeff{P_{4k+1}}{4k-2}&a&-2k\\
      \coeff{P_{4k+1}}{4k-3}&e&k\\
      \coeff{P_{4k+1}}{4k-4}&\rho&k
    \end{array}
    \]
    In the fourth and fifth rows the possible boundary contributions from
    $A,B$ vanish when $k=1$ and, when $k\ge2$, are parameter-free and hence
    subtractable by \Cref{lem:Rk2l-top-boundary}.  By
    $(4k+1)$-admissibility, every displayed
    slope is invertible.  Thus $H_2,H_4,\tilde H_4$ and all
    five outer parameters are derivable from $P_{4k+1}$.

    We now check the cutoffs directly.  The coefficients of $U$ above
    degree $4(k-1)$ and of $V$ above the same boundary are coefficients of
    the known powers.  At the boundary, the parameter-free top-boundary
    assertion \Cref{lem:Rk2l-top-boundary} supplies the one remaining
    correction.  Descending from the top five rows therefore derives
    $U_j$ from rows of $P_{4k+1}$ at least $j+1$ and $V_j$ from rows at least
    $j$, for all $j\ge4(k-1)$.  After the five outer parameters and the
    powers have been substituted, the residual
    \[
      x(U-H_4^k)+\bigl(V-\tilde H_4^k\bigr)=xA+B
    \]
    is coefficient-triangular on rows below $4(k-1)$ by
    \Cref{lem:Rk2l}(3).  Every outer parameter, and hence every coefficient
    of these powers, was recovered from rows at least $4(k-1)$.  For an inner
    row $j<4(k-1)$ this is at least $j+1$.  Substituting the top-block formulas
    into the triangular decoder therefore preserves both cutoffs: residual
    rows at least $j+1$ recover $A_j$, and residual rows at least $j$ recover
    $B_j$.  Hence $A,B$, and
    therefore $U,V$, satisfy the two compatibility cutoffs on the whole
    window.  If $k\ge2$, the same argument extracts
    $\alpha_0,\ldots,\alpha_{4k-5}$; for $k=1$ that internal block is empty.
    This proves the decoder and compatibility claims.  For the realization,
    the quadratic $H_2$ costs one product, and the shared
    $T_{2k,2}$ call costs $\mu(2k,1)=2k-1$ products by
    \Cref{lem:T-multiplication-count}.  The same call outputs both components
    and forms $H_4$ as an intermediate wire.  Thus the single shared circuit
    has $2k$ products and records both $H_2$ and $H_4$.
\end{proof}

\begin{lemma}[A decodable $Q_{4k+1}$ given $H_2$]\label{lem:Q4k+1-from-H2}
    Let $k\ge 1$, assume that $\mathbb F$ is $(4k+1)$-admissible, and let
    $H_2\in\mathbb{F}[x]$ be monic of degree $2$.

    Define
    \[
      \hat H_2 \coloneqq H_2+\alpha_{4k-1},
      \qquad
      \tilde H_2 \coloneqq \hat H_2+\alpha_{4k-2},
    \]
    and let
    \[
      S^{(1)}_2 \coloneqq T^{(1)}_{2k,2}[\alpha_0,\ldots,\alpha_{4k-3}](x,\hat H_2),
      \qquad
      S^{(2)}_2 \coloneqq T^{(2)}_{2k,2}[\alpha_0,\ldots,\alpha_{4k-3}](x,\hat H_2,\tilde H_2).
    \]
    Finally define
    \[
      Q_{4k+1}[\alpha_0,\ldots,\alpha_{4k}](x,H_2)\coloneqq (x+\alpha_{4k})\,S^{(1)}_2(x)+S^{(2)}_2(x).
    \]
    Then $Q_{4k+1}$ is decodable given $H_2$, and the polynomial $H_4$ computed in the internal even-$k$, $l=1$ step is derivable from $Q_{4k+1}$ given $H_2$.
\end{lemma}
\begin{proof}
    Write $H_2=x^2+bx+c_0$ and put
    \[
      \gamma=\alpha_{4k-1},\quad \rho=\alpha_{4k-2},\quad
      a=\alpha_{4k-3},\quad e=\alpha_{4k-4},\quad
      \beta=\alpha_{4k}.
    \]
    The shared base computes
    \[
      H=H_2+\gamma,\qquad
      H_4=H^2-(x+a)^2+e,\qquad
      \tilde H_4=H_4+\rho,
    \]
    and then
    \[
      (S^{(1)}_2,S^{(2)}_2)
       =T_{k,4}[\alpha_0,\ldots,\alpha_{4k-5}]
          (x,H,H_4,\tilde H_4).
    \]
    Since $\tilde H_4-H_4=\rho$ and the field is $4k$-admissible, the
    hypotheses of \Cref{lem:Rk2l} hold for this $T_{k,4}$ call.
    As in the preceding proof, the top-boundary lemma and a direct Cauchy-product
    expansion give the descending affine pivots
    \[
    \begin{array}{c|c|c}
      \coeff{Q_{4k+1}}{4k}&\beta&1\\
      \coeff{Q_{4k+1}}{4k-1}&\gamma&2k\\
      \coeff{Q_{4k+1}}{4k-2}&a&-2k\\
      \coeff{Q_{4k+1}}{4k-3}&e&k\\
      \coeff{Q_{4k+1}}{4k-4}&\rho&k
    \end{array}
    \]
    Every entry not containing the displayed new parameter is a polynomial in
    the coefficients of the given $H_2$ and in parameters recovered in earlier
    rows.  Hence all five parameters, and therefore $H,H_4,\tilde H_4$, are
    derivable.

    Given $H,H_4,\tilde H_4$, \Cref{lem:Rk2l}(3) says directly that
    $(S^{(1)}_2,S^{(2)}_2)$ is compatible on its remainder window.  The five
    top rows have already recovered those three polynomials and $\beta$, so
    the known-shift clause of \Cref{lem:x-alpha-extraction} applies to
    $Q_{4k+1}=(x+\beta)S^{(1)}_2+S^{(2)}_2$ and derives the pair.  This clause
    permits the full top row in the compatibility window; it does not try to
    recover $\beta$ from the generic top-row formula.
    If $k\ge2$, subtracting $H_4^k$ and $\tilde H_4^k$ and applying
    \Cref{lem:Rk2l}(3) at $(k,2)$ extracts
    $\alpha_0,\ldots,\alpha_{4k-5}$; for $k=1$ this block is empty.  Thus
    $Q_{4k+1}$ is decodable
    given $H_2$, and the already recovered $H_4$ is the required byproduct.
\end{proof}
